\chapter{The Object Map}

APFS resolves virtual object identifiers to physical block
addresses through a B-tree called an object map (OMAP). There is
one per container (resolving volume superblocks and ephemerals)
and one per volume (resolving the FS-tree root and its
descendants). The OMAP is the mechanism by which a copy-on-write
filesystem maintains stable object identity across mutations.

The lookup takes \emph{two} keys, and the second is non-negotiable.

\section{Lookup By (OID, Max-XID)}

\textbf{Spec.} An OMAP key is \texttt{(oid, xid)}. An OMAP value names
a physical address, a size, and a flag word. The OMAP can hold
multiple entries for a single OID, each at a different XID,
corresponding to different versions of the object over time.

\textbf{Spec.} An OMAP lookup specifies an \texttt{oid} and a
\emph{max XID}. The B-tree returns the entry with the highest stored
XID that is \emph{less than or equal to} the requested max XID. The
semantics are lower-bound rather than exact match: the lookup is
asking "give me the version of \texttt{oid} that was current at or
before this scan."

The resolver input is a triple: the OMAP context (which OMAP owns
the OID), the OID, and the scan XID. The OID by itself does not
determine which version of the object is meant; the scan XID does.

The scan XID is fixed for the entire run --- the one selected from
the descriptor area in chapter~3. Mixing scan XIDs across
resolutions reintroduces the incoherence that pinning the
checkpoint was supposed to prevent.

\section{Hard Stops At OMAP Open}

When the parser opens an OMAP --- reads the
\texttt{omap\_phys\_t} root object and validates it --- the
\texttt{omap\_phys.flags} field gates whether the OMAP can be used at
all.

\begin{description}
  \item[\texttt{OMAP\_PHYS\_ENCRYPTING}.]
  \item[\texttt{OMAP\_PHYS\_DECRYPTING}.]
  \item[\texttt{OMAP\_PHYS\_KEYROLLING}.]
  \item[\texttt{OMAP\_PHYS\_CRYPTO\_GENERATION\_FLAG}.]
\end{description}

Any of these set means an encryption operation is in flight on the
container or the volume. Raw mode cannot be supported in such a
state: the bytes the reader sees do not yet match the bytes the
running system will commit. The parser fails closed; the caller
falls back to a supported API.

\textbf{Observation.} Any \emph{unknown} bit in \texttt{omap\_phys.flags}
is also a hard stop. APFS reserves bits over time; an unknown bit
means the running OS knows something the parser does not, and the
parser must not infer behaviour. This is the format-drift discipline
of chapter~11 applied at the smallest object that hosts it.

\section{Hard Stops At OMAP Lookup}

Each value entry returned by an OMAP lookup carries its own flag
word. The flags decide whether the value is usable.

\begin{description}
  \item[\texttt{OMAP\_VAL\_DELETED}.] Negative result, not an error.
    The OID had a mapping at some earlier XID but is no longer
    visible at the scan XID. The caller interprets this as
    "the object is not present" and proceeds accordingly. The
    parser does not hard-stop.
  \item[\texttt{OMAP\_VAL\_ENCRYPTED}.] The value is per-object
    encrypted. Raw mode cannot decrypt; hard stop.
  \item[\texttt{OMAP\_VAL\_NOHEADER}.] The target block does not
    carry an \texttt{obj\_phys\_t} header at all. A structural
    assumption the rest of the parser depends on has been broken;
    hard stop.
  \item[\texttt{OMAP\_VAL\_CRYPTO\_GENERATION}.] Encryption-generation
    mismatch between the OMAP entry and the volume; hard stop.
  \item[unknown bits.] Format drift; hard stop.
\end{description}

The discipline is tighter at lookup time than at open time only
because lookup is also where the parser sees per-entry encryption
behaviour. The principle is the same: known flag means known
behaviour; unknown flag means the parser does not understand the
running system and refuses to guess.

\section{Empirical Correction: B-Tree TOC Layout (EX-12)}

\textbf{Observation.} The B-tree leaf TOC \texttt{v\_off} field
measures the distance from the \emph{end} of the value area to the
\emph{start} of the value. The value runs forward from that start
for \texttt{fixed\_value\_size} bytes (or \texttt{v\_len} for
variable-length values).

The reference admits a reading under which the value extends
\emph{backward} from \texttt{v\_off}. That reading produces
nonsense OMAP values --- \texttt{flags=0x10ffff} on the proof
fixture --- and corrupts every downstream lookup. The forward
reading produces clean values (\texttt{flags=0x0},
\texttt{size=4096}, \texttt{paddr=...}) and matches every
converged reader. A synthetic OMAP regression test catches any
reintroduction of the bug.

\section{Cross-Tool Validation}

EX-12 paired the project's parser with \texttt{go-apfs identitydump}
against the same \texttt{/dev/rdiskN}. The probe:

\begin{enumerate}
  \item built a fresh proof fixture,
  \item ran the project's parser, recording every resolved
    \texttt{(oid, paddr, xid)} sample published by the OMAP,
  \item replayed object-header validation rules in Python at each
    published \texttt{paddr},
  \item ran \texttt{go-apfs identitydump} against the same raw
    container,
  \item compared the two tools' agreement on the FS-tree root
    \texttt{oid}.
\end{enumerate}

The two tools agreed on the FS-tree root OID. They disagreed on
which scan XID to select: this project's parser picked the highest
visible XID; \texttt{go-apfs} picked one lower. The divergence is
recorded as a caveat
(\texttt{go\_apfs\_active\_state\_observation}); both tools' OMAP
lookups internally use a consistent scan XID and produce internally
consistent root resolutions, so the cross-tool oracle survives.

Two tools that disagree on selection but agree on the OID reached
through that selection remain useful as oracles for each other,
provided the selection difference is recorded.

\section{What The OMAP Does Not Do}

The OMAP does not resolve objects \emph{within} the FS tree.
FS-tree internal-node entries also carry OIDs, and those OIDs
require OMAP resolution at the scan XID --- through the
\emph{volume} OMAP, not the container OMAP. The next chapter
develops this rule and the fixture that surfaced it.
